\chapter{Введение}

\section{Цель}  

Сформировать списки контроля доступа. Назначить списки
на интерфейсы. Проверить работоспособность сети и сетевых сервисов до и после
назначения списков на интерфейс
\section{Задачи}

\begin{enumerate}
\item Соединить устройства в сеть
\item Настроить IP - адресацию на интерфейсах и имена устройств согласно
таблице
\item Включить telnet и ssh сервер на маршрутизаторе R1.
\item Создать списки контроля доступа Test1, Test2.
\item В списке контроля доступа Test1 настроить: запретить входящий трафик по
telnet, разрешить SSH и ICMP.
\item В списке контроля доступа Test2 настроить: запретить входящий трафик по
telnet, SSH и ICMP.
\item На интерфейсе Gi1/0/3 не ставить никаких правил. Все разрешить.
Проверить работу сети. Затем в списке контроля доступа Test2, разрешить
трафик по протоколу ICMP
\end{enumerate}
 
\begin{figure}[H]  
    \centering
    \includegraphics[width=0.75\textwidth]{top.png}
    \caption{Пример собранной сети}
\end{figure}

\begin{figure}[H] 
    \centering
    \includegraphics[width=0.75\textwidth]{table.png}
    \caption{Таблица с настройками}
\end{figure}

\endinput